\lesson{1}{Sep 07 2021 Tue 11:09}{The Science of Chemistry}{Unit 1}

\begin{definition}[Chemistry]
	Chemistry is the study of composition and the structure of materials and the changes they undergo.

	This includes studying materials found in our:

	\begin{itemize}
		\item Ocean
		\item Atmosphere
		\item Environment
		\item Underground
		\item Etc.\ldots
	\end{itemize}

	Chemistry is one of many fields of science that investigate phenomena such as climate change to discover their causes.
\end{definition}

\begin{definition}[Science]
	Science must be based on empirical observations, experimentation, explanations based on logical reasoning.
	\begin{itemize}
		\item OBSERVABLE: Science attempts to explain natural phenomena by analyzing and observing the world and testing ideas about it.
		\item TESTABLE: Science must be able to answer a testable question using observation and experimentation. Investigations must produce empirical evidence that can be observed or measured to be considered science.
		\item REPLICABLE: Empirical evidence can be replicated, or reproduced, and verified by other scientists if they conduct the same tests and under the same conditions.
		\item RELIABLE: The more an experiment is repeated, with the same outcomes, the more reliable the evidence becomes. Evidence with bias also increases its reliability.
		\item FLEXIBLE: Science is an ever-changing body of knowledge as new observations are made through experimentation. As new information is discovered, new evidence can add to current evidence, allowing scientists to improve their theories.
	\end{itemize}
\end{definition}

Many fields of knowledge, such as philosophy and art add to our view of the world. They help us appreciate the beauty of the world, interact with one another and decide what is wrong and right. But, since there aren't any observations and tests being applied in those fields and everything in them are just beliefs and opinions, then it cannot be called science.

These two questions will help you determine if the question can be answered with science:

\begin{itemize}
	\item If the question is asking about an opinion or a moral value, it's not something that can be measured using scientific process. Therefor, it cannot be answered with science.
	\item If the answer to the question cannot be tested and observed, it is not considered science.
\end{itemize}

\subsubsection*{The Scientific Method}

A scientific method is a series of steps for investigating questions and testing ideas. There are several versions of the \textbf{scientific method}, but all of the versions are based on rational thinking, inquiry, and experimentation. The \textbf{scientific method} includes five main steps:

\begin{itemize}
	\item Question: A scientific method always starts with a question. As long as it's testable, you can use science to find the answer.
	\item Research: It's important that you always explore what other scientists before have observed and discovered because that information may be able to solve your current question.
	\item Hypothesis: Most hypotheses are made with the \textbf{"if/then"} statements. If "this" happens, then "that" will take place. This helps you know how one thing can affect another and gives you variables to test.
	\item Testing: An experiment allows you to test your hypothesis to determine if you're correct or incorrect.
		\item Independent Variable: This variable is the factor the scientist has chosen to change in the experiment. A good experiment should change only one independent variable because this allows the scientist to focus only on the outcome of that specific variable.
		\item Dependent Variable: This variable is the factor that changes in response to the independent variable in an experiment.
		\item Controlled Variables: These variables are the factors the scientist chooses to keep constant over the course of the experiment. By keeping all other factors constant, anything happens to the dependent variable is caused by the independent variable.
\end{itemize}

\newpage
